% Generated draft: requires mathematical and source review.
\documentclass[11pt]{article}
\usepackage[margin=25mm]{geometry}
\usepackage{fontspec}
\setmainfont{Latin Modern Roman}
\usepackage{amsmath,amssymb,mathtools,unicode-math}
\setmathfont{Latin Modern Math}
\usepackage{enumitem,needspace,xurl,hyperref}
\setlength{\parindent}{0pt}
\setlength{\parskip}{5pt}
\setlength{\emergencystretch}{3em}
\newcommand{\N}{\mathbb{N}}
\newcommand{\Z}{\mathbb{Z}}
\newcommand{\Q}{\mathbb{Q}}
\newcommand{\R}{\mathbb{R}}
\newcommand{\C}{\mathbb{C}}
\newcommand{\F}{\mathbb{F}}
\newcommand{\A}{\mathbb{A}}
\newcommand{\PP}{\mathbb{P}}
\newcommand{\E}{\mathbb{E}}
\newcommand{\eps}{\varepsilon}
\newcommand{\dd}{\,\mathrm{d}}
\begin{document}
\section*{UnsolvedMath ID: 30001706}
\textbf{Source problem code:} catalogue entry
\subsection{Definitions and mathematical statement}
Let \(G\) be a connected Lie group with Lie algebra \(\mathfrak g\), let \(K\subseteq G\) be a compact subgroup with Lie algebra \(\mathfrak k\), and let \(O=\operatorname{Ad}^{*}(G)\lambda\subseteq\mathfrak g^{*}\) be a coadjoint orbit, where \(\lambda\in\mathfrak g^{*}\). The orbit projection is the restriction map
\[
\mu_K:O\longrightarrow\mathfrak k^{*},\qquad
\mu_K(\xi)=\xi|_{\mathfrak k}.
\]
With the compatible sign convention for the Kirillov--Kostant--Souriau symplectic form, this is a moment map for the restricted coadjoint action of \(K\).

To describe the usual chamber construction, let \(K^{\circ}\) be the identity component of \(K\), choose a maximal torus \(T\subseteq K^{\circ}\), and choose a closed positive Weyl chamber \(\mathfrak t^{*}_{+}\). A \(K\)-invariant inner product on \(\mathfrak k\) allows this chamber to be regarded as a subset of \(\mathfrak k^{*}\). Define the chamber image and its closure by
\[
S_K(O)=\mu_K(O)\cap\mathfrak t^{*}_{+},\qquad
\overline{S_K(O)}=\operatorname{cl}_{\mathfrak t^{*}_{+}}S_K(O).
\]
These definitions do not assert that either set is convex or polyhedral.

\textbf{Source problem.} Given such \(G\), \(K\), and \(O\), determine explicit equations describing the moment polyhedron \(\Delta_K(O)\) of the orbit projection \(\mu_K\).

\textbf{Unresolved qualifications.} The supplied statement does not define \(\Delta_K(O)\), specify whether it denotes the chamber image or its closure, or supply hypotheses guaranteeing convexity or polyhedrality. It also does not specify whether ``equations'' includes affine inequalities or requires a finite description. The chamber construction above uses \(K^{\circ}\); if disconnected compact subgroups are intended, the treatment of their additional component action must be specified. No compactness of \(G\) or \(O\), properness of \(\mu_K\), reductivity, or holomorphic-orbit hypothesis is imposed in the supplied statement. Accordingly, the requested description is recorded in its supplied generality, with these conventions and hypotheses unresolved.
\subsection{Short English statement}
For a coadjoint orbit of a connected Lie group, describe explicitly the moment image obtained by restricting covectors to the Lie algebra of a compact subgroup and passing to a positive Weyl chamber. The source calls the target a moment polyhedron, but leaves its closure convention and the hypotheses ensuring polyhedrality unspecified. Its partially solved label is supplied dated curation and requires independent verification.
\subsection{Sources}
\begin{itemize}
\item Supplied local problem registry record, permanent UnsolvedMath ID \(30001706\), title \emph{Moment Polyhedra of Noncompact Coadjoint Orbits}, exact source problem code \texttt{catalogue entry}; collection: \emph{Oberwolfach Reports Open Problems}. The record gives proposed year 2011 and names no individual proposer.
\item Supplied report citation: \emph{Geometric Quantization in the Non-compact Setting} (2011), Oberwolfach Reports, pp.~425--521, DOI \texttt{10.4171/owr/2011/09}. Original report: \url{https://doi.org/10.4171/owr/2011/09}. The supplied background describes Oberwolfach Reports as records of mathematical workshops and their open-problem discussions.
\item Supplied literature source: Deltour, \emph{Kirwan polyhedron of holomorphic coadjoint orbits}, \url{https://arxiv.org/abs/1104.1909}.
\item Literature assessment dated exactly \texttt{2026-08-21}: the supplied curation labels the problem partially solved and states that Berenstein--Sjamaar solve the compact case, while Deltour gives affine equations for holomorphic noncompact coadjoint orbits; it states that the full extracted generality remains broader. This is supplied dated curation, not an independently verified literature assessment. No separate bibliographic reference for the Berenstein--Sjamaar claim is supplied.
\item The registry labels the statement \texttt{corrected\_verified} and reports a manual semantic, typographic, and LaTeX review of its cleanup. These are supplied editorial assertions, not mathematical verification of this draft or independent inspection of the cited sources.
\end{itemize}
\end{document}
